% Generated draft: requires mathematical and source review.
\documentclass[11pt]{article}
\usepackage[margin=25mm]{geometry}
\usepackage{fontspec}
\setmainfont{Latin Modern Roman}
\usepackage{amsmath,amssymb,mathtools,unicode-math}
\setmathfont{Latin Modern Math}
\usepackage{enumitem,needspace,xurl,hyperref}
\setlength{\parindent}{0pt}
\setlength{\parskip}{5pt}
\setlength{\emergencystretch}{3em}
\newcommand{\N}{\mathbb{N}}
\newcommand{\Z}{\mathbb{Z}}
\newcommand{\Q}{\mathbb{Q}}
\newcommand{\R}{\mathbb{R}}
\newcommand{\C}{\mathbb{C}}
\newcommand{\F}{\mathbb{F}}
\newcommand{\A}{\mathbb{A}}
\newcommand{\PP}{\mathbb{P}}
\newcommand{\E}{\mathbb{E}}
\newcommand{\eps}{\varepsilon}
\newcommand{\dd}{\,\mathrm{d}}
\begin{document}
\section*{UnsolvedMath ID: 30004955}
\textbf{Source problem code:} catalogue entry
\subsection{Definitions and mathematical statement}
\textbf{Unrecoverable mathematical statement.} The supplied record contains workshop introductory prose, with no mathematical problem, definitions, hypotheses, domains, boundary conditions, quantifiers, or requested conclusion from which a precise statement can be recovered. Its title, ``Unspecified Structural Graph Theory Problem,'' does not specify a mathematical target.

The abstract describes a graph theory workshop intended to bring together researchers at various career stages to discuss developments across graph theory and its connections with mathematics and computer science. It mentions the Product Structure Theorem and progress towards the resolution of Hadwiger's Conjecture, and describes talks and collaboration opportunities despite disruption by the COVID pandemic. Neither named topic is formulated as the problem attached to this record; selecting either would introduce an unsupported mathematical target.

The supplied statement status is \texttt{unrecoverable}. The supplied statement-verification note characterizes the passage as workshop-level introductory prose rather than a standalone mathematical problem and says that no separate clean statement is promoted. The supplied \texttt{open} label reflects a curation uncertainty rule for an abstract with no specific claim to verify; it does not establish that an identifiable mathematical problem is open.
\subsection{Short English statement}
The supplied workshop abstract does not specify a mathematical problem. Its references to the Product Structure Theorem and Hadwiger's Conjecture do not identify an intended statement. A statement-only draft must retain this unresolved status pending human review of the provenance and intended target.
\subsection{Sources}
\begin{itemize}
\item Local UnsolvedMath registry: \texttt{data/unsolvedmath/problems.json}, permanent ID \texttt{30004955}, source code \texttt{catalogue entry}, title ``Unspecified Structural Graph Theory Problem,'' in the collection ``Oberwolfach Reports Open Problems.'' The supplied record retains the original passage for provenance and supplies no clean statement or named proposer.
\item Supplied source citation: \emph{Graph Theory (2023)}, Oberwolfach Reports, DOI 10.4171/owr/2022/1, pp.\ 5--78. Supplied source and background link: \url{https://doi.org/10.4171/owr/2022/1}. The proposed-year field and year tag give 2022; the citation title gives 2023. Their relationship is not explained in the supplied record.
\item Supplied literature source, labelled ``original OWR report'': \url{https://doi.org/10.4171/OWR/2022/1}. Literature assessment checked at \texttt{2026-08-22}: supplied dated curation, not independently verified here. The assessment labels the record open under an uncertainty rule because the extracted workshop abstract contains no specific mathematical problem or claim to verify. The supplied background asks that time-sensitive status and cited claims be verified independently.
\end{itemize}
\end{document}
